\documentclass[11pt]{report}

\usepackage[utf8]{inputenc}
\usepackage[T1]{fontenc}
\usepackage{amsmath}
\usepackage{booktabs}
\usepackage[margin=1in]{geometry}
\usepackage{hyperref}

\title{Report Title}
\author{Your Name \\ \small Your Organisation}
\date{\today}

\begin{document}

\maketitle

\tableofcontents

\chapter{Introduction}
\label{ch:intro}

This report template has a title page, a table of contents, two chapters,
a figure, a table and a bibliography. Replace this text with your own.

\section{Background}

State the question the report answers and why it matters. Earlier work
can be cited from the bibliography file~\cite{lamport1994}.

\section{Scope}

The report covers:
\begin{itemize}
  \item the first topic;
  \item the second topic;
  \item the limits of this work.
\end{itemize}

\chapter{Findings}
\label{ch:findings}

\section{Overview}

Chapter~\ref{ch:intro} set out the scope. This chapter presents the
results. A key relation is
\begin{equation}
  y = a x + b .
  \label{eq:line}
\end{equation}

\begin{figure}[h]
  \centering
  \fbox{\parbox[c][4cm][c]{0.6\textwidth}{\centering Figure placeholder}}
  \caption{A placeholder for a figure.}
  \label{fig:placeholder}
\end{figure}

\section{Results}

Table~\ref{tab:results} lists example values.

\begin{table}[h]
  \centering
  \begin{tabular}{lrr}
    \toprule
    Item & Value & Change \\
    \midrule
    First & 10 & +2 \\
    Second & 25 & -1 \\
    Third & 40 & +5 \\
    \bottomrule
  \end{tabular}
  \caption{Example values.}
  \label{tab:results}
\end{table}

\section{Conclusion}

Summarise the findings and the next steps here.

\bibliographystyle{plain}
\bibliography{references}

\end{document}
